% META: {"id": "TNSI-ARCH-EX-013", "chapitre": "TNSI-ARCHITECTURES-MATERIELLES-SY", "type_objet": "exercice", "capacites": ["TNSI-ARCHITECTURES-MATERIELLES-SY-C3"], "capacites_codes": ["C3"], "methodes": ["M1"], "parcours": 1, "competences": ["programmer"], "duree_min": 8, "mode_creation": "ex_nihilo", "sources_inspiration": [], "fichier_tex": "chapitres/TNSI-ARCHITECTURES-MATERIELLES-SY/exercices/TNSI-ARCH-EX-013.tex", "corrige_tex": "chapitres/TNSI-ARCHITECTURES-MATERIELLES-SY/corriges/TNSI-ARCH-CO-013.tex", "status": "approved"}
% BEGIN-VERIFY
% def tourniquet(processus, tranche, duree_totale):
%     ordre_execution = []
%     file = list(processus)
%     temps_restant = dict(duree_totale)
%     while file:
%         p = file.pop(0)
%         execute = min(tranche, temps_restant[p])
%         ordre_execution.append((p, execute))
%         temps_restant[p] -= execute
%         if temps_restant[p] > 0:
%             file.append(p)
%     return ordre_execution
% resultat = tourniquet(["X", "Y"], 3, {"X": 4, "Y": 6})
% assert resultat == [("X", 3), ("Y", 3), ("X", 1), ("Y", 3)]
% END-VERIFY
\begin{exercice}{TNSI-ARCH-EX-001}{1}{15}
En utilisant la fonction \lstinline{tourniquet} du cours, déterminer l'ordre d'exécution pour deux processus \texttt{X} (durée totale $4$) et \texttt{Y} (durée totale $6$), avec une tranche de temps de $3$. Combien de tranches de temps sont-elles nécessaires au total pour que les deux processus se terminent ?
\end{exercice}
